%!TEX root = ../main.tex
\chapter{Confronto}
\label{cap:confronto}

Questo capitolo descrive il protocollo sperimentale di questa tesi e ne
riporta i risultati. Definisce anzitutto le metriche di valutazione
(§\ref{sec:metriche}) e il protocollo sperimentale
(§\ref{sec:protocollo}); presenta poi il confronto di base tra le
baseline classiche e MetaSTGAT Pro insieme ai tre studi sperimentali di
questa tesi, sul peso $\alpha$ della componente di equità, sull'ablation
della procedura di addestramento e sui meccanismi di comunicazione
spaziale alternativi (§\ref{sec:confronti}), e ne riporta i risultati
(§\ref{sec:risultati-baseline}).

\section{Metriche di valutazione}
\label{sec:metriche}

Nella letteratura sul Traffic Signal Control le metriche di valutazione
proposte sono numerose.
Questo capitolo introduce le 3 metrice principali utilizzate, coerentemente con il duplice obiettivo di
questo lavoro, ovvero efficienza ed equità (come motivato nel Capitolo~\ref{cap:introduzione}):
il \emph{travel time} (tempo di viaggio) e il \emph{trip completion} (veicoli che
completano il percorso) per l'efficienza, il \emph{max wait} (l'attesa
massima registrata da un veicolo ad un incrocio) per l'equità.
C'è da sottolineare che né \emph{travel time} né
\emph{trip completion} sono metriche che appaiono esplicitamente nella
ricompensa dei modelli. Pertanto qualunque risultato su di esse sarà
l'indiretta conseguenza di come la reward ha indirizzato l'apprendimento
della policy.

\paragraph{Travel time.} Per ogni veicolo $v$, il tempo di viaggio è
\begin{equation}
  tt_v =
  \begin{cases}
    t_{\text{arrivo}}(v) - t_{\text{spawn}}(v)     & \text{se } v \text{ è arrivato entro la fine dell'episodio} \\
    t_{\text{fine episodio}} - t_{\text{spawn}}(v) & \text{altrimenti (un lower bound del vero valore)}
  \end{cases}
\end{equation}
Una delle due metriche cardini in questo studio è il travel time medio $\overline{tt}$,
ovvero la media di $ tt_v $ calcolata sull'intera popolazione di veicoli generati nell'episodio.
Tale metrica non considera soltanto i veicoli arrivati, ma include anche i
veicoli che non hanno completato il percorso entro la fine della simulazione.
La ragione è evitare un bias di sopravvivenza: un
modello che congestiona la rete e lascia arrivare solo una minoranza di
veicoli su percorsi ottimizzati otterrebbe, calcolando la media solo sui
veicoli arrivati, un risultato positivo ma non rappresentativo del comportamento reale della rete.

La seconda vera metrica cardine in questo lavoro di tesi è il travel time massimo $\overline{tt}_{\max}$,
ovvero il massimo tra tutti i travel time dei veicoli generati, anche in questo caso indipendentemente
dal completamento del percorso, per motivi analoghi a quelli descritti per il travel time medio.

\paragraph{Trip completion.} Questa metrica misura la percentuale di veicoli che hanno completato il percorso entro la
fine dell'episodio, escludendo dal totale i veicoli che in assenza di semafori e in condizioni di flusso libero non riuscirebbero comunque
a completare il proprio percorso.
In questo studio questa metrica viene indicata con $tc$. Il dettaglio di
$tc$ per ogni singolo modello e ogni singola configurazione è in
Tabella~\ref{tab:tc-adjusted}, a fine capitolo.

\paragraph{Max wait.} Questa metrica riguarda invece il massimo periodo di tempo
per cui un veicolo è rimasto bloccato allo stesso incrocio.
Dunque questo valore non si azzera nel momento in cui un veicolo avanza nella corsia, ma si annulla
soltanto se la vettura supera l'incrocio; ciò significa che max wait non considera solamente
il lasso di tempo in cui il veicolo è rimasto fermo, ma tiene conto anche del tempo in cui lo stesso si trova in movimento all'interno della corsia senza però oltrepassare l'intersezione.
Essa viene calcolata prendendo il massimo valore tra tutti i veicoli e
tutte le intersezioni.

Come per il travel time, include anche gli stalli ancora in corso al
termine dell'episodio.
In questo studio max wait viene indicato con $\text{wait}_{\max}$.

La presente metrica non va confusa con $\mathrm{MaxRedWait}_i^t$, il termine
usato nella ricompensa (Capitolo~\ref{cap:tsc-modelli}) per
penalizzare l'attesa durante il training. $\mathrm{MaxRedWait}_i^t$ è
calcolato a ogni singolo passo di decisione, su una sola intersezione alla
volta, ma soprattutto è limitata al campo visivo dell'incrocio.
$\text{wait}_{\max}$ è invece calcolato una sola volta a fine
episodio, come massimo su tutta la rete, e fa riferimento all'intera corsia.

\section{Protocollo sperimentale}
\label{sec:protocollo}

Ogni modello RL è addestrato sulle tre configurazioni di training presentate nella Tabella~\ref{tab:split}, riportata
per comodità di lettura qui sotto.

\begin{table}[htbp]
  \centering
  \footnotesize
  \begin{tabular}{@{}lllll@{}}
    \toprule
    Griglia    & Spaziatura & Profilo di traffico       & Utilizzo    & Configurazione                          \\
    \midrule
    $4\times4$ & 100m       & Complete (arteria riga 0) & Train, Test & \texttt{config\_4x4\_100m\_train1}      \\
    $4\times4$ & 100m       & Complete (arteria riga 1) & Train, Test & \texttt{config\_4x4\_100m\_train2}      \\
    $4\times4$ & 100m       & Complete (arteria riga 2) & Train, Test & \texttt{config\_4x4\_100m\_train3}      \\
    $4\times4$ & 100m       & Flat                      & Validation  & \texttt{config\_4x4\_100m\_6k\_flat}    \\
    $4\times4$ & 100m       & Peak                      & Test        & \texttt{config\_4x4\_100m\_6k\_peak}    \\
    $4\times4$ & 200m       & Flat                      & Test        & \texttt{config\_4x4\_200m\_6k\_flat}    \\
    $4\times4$ & 200m       & Peak                      & Test        & \texttt{config\_4x4\_200m\_6k\_peak}    \\
    $5\times5$ & 100m       & Flat                      & Test        & \texttt{config\_5x5\_100m\_9.4k\_flat}  \\
    $5\times5$ & 100m       & Peak                      & Test        & \texttt{config\_5x5\_100m\_9.4k\_peak}  \\
    $6\times6$ & 100m       & Flat                      & Test        & \texttt{config\_6x6\_100m\_11.5k\_flat} \\
    $6\times6$ & 100m       & Peak                      & Test        & \texttt{config\_6x6\_100m\_11.5k\_peak} \\
    \bottomrule
  \end{tabular}
\end{table}

Seguendo un ciclo fisso, le tre configurazioni si alternano episodio per
episodio. A fine training, l'ultimo checkpoint e quello che ha ottenuto il
valore di loss più basso durante l'intero addestramento sono confrontati
sulla configurazione di validazione \texttt{config\_4x4\_100m\_6k\_flat},
mediando $\overline{tt}$ su 5 seed diversi.
Questa fase serve a evitare due problemi che si concatenano tra loro.
Il primo è che in un ciclo di
apprendimento online la politica non converge necessariamente in modo
monotono, dunque gli ultimi episodi addestrati potrebbero allontanare la policy da un
punto di minimo quantomeno locale.
Il secondo è che la scelta del modello che ha registrato la loss minima non garantisce
le capacità di generalizzazione del modello, in quanto una loss bassa potrebbe
anche essere sintomo di una policy che ha soltanto memorizzato delle condizioni particolari delle
configurazioni di training, senza però aver realmente appreso una politica in grado di generalizzare.

Il modello che risulta essere migliore in questa fase di selezione, viene poi utilizzato nella fase di test.
Ogni valutazione, sia in validation sia nella batteria di test che descriveremo ora, avviene con $\varepsilon=0$:
nessuna esplorazione residua, solo la politica appresa allo stato puro. Inoltre i risultati ottenuti in ciascuna
configurazione testata, sono mediati su 5 seed diversi, analogamente a quanto fatto per $\overline{tt}$ nella fase di validazione.

Misurare la capacità di generalizzazione dei modelli è uno degli obiettivi
dichiarati di questo lavoro (Capitolo~\ref{cap:introduzione}), un asse che
gran parte della letteratura sul controllo semaforico appreso non copre
sistematicamente, in quanto la prassi più diffusa, seguita anche dall'articolo
originale di MetaSTGAT \citep{wang2022metastgat}, è addestrare e valutare un modello sulla stessa
rete stradale, senza mai trasferirlo a una topologia diversa.

Per questo motivo la batteria di test di questo lavoro include sia le 3
configurazioni di training sia 7 configurazioni aggiuntive mai viste
durante l'addestramento. Le prime permettono un confronto diretto con la
prassi di letteratura appena descritta, riportando un risultato ottenuto
nelle stesse condizioni con cui la maggior parte degli altri lavori
riporta il proprio; le seconde, assenti nella maggior parte degli studi
citati, misurano invece se le prestazioni osservate si confermano anche con il cambio
di topologia e di condizioni di traffico.

Ogni modello, RL o baseline, è dunque valutato sistematicamente sulle
stesse 10 configurazioni: le 3 varianti di training stesse (per misurare la
prestazione in-distribuzione) più 7 configurazioni mai viste in training
(4x4 a 200m invece di 100m, $5\times5$ e $6\times6$ a densità sia
\emph{flat} sia \emph{peak}), usate per misurare la generalizzazione a
topologie e densità di traffico diverse.

\section{Confronti proposti}
\label{sec:confronti}

Il confronto sperimentale di questo capitolo è organizzato in tre studi indipendenti.
Ognuno di essi isola un solo asse di variazione rispetto al
modello di riferimento ($\alpha=0.08$, MetaSTGAT a un layer) a cui ci riferiremo con il termine \emph{MetaSTGAT Pro},
lasciando invariato tutto il resto. Verranno analizzati nel seguito i tre studi proposti.
Ciascun confronto conterrà al suo interno, come riferimento, le due baseline classiche, Fixed-Time e MaxPressure, e il modello \emph{MetaSTGAT Pro}.

\subsection{Studio del peso \texorpdfstring{$\alpha$}{alpha}}

Il primo studio (§\ref{sec:risultati-baseline}) pone a confronto tre varianti del
modello MetaSTGAT Pro che differiscono solo per il valore del peso $\alpha$ riguardante la penalità anti-starvation
nella ricompensa (Eq.~\ref{eq:reward-custom}).
In particolare vengono valutati i valori di $\alpha \in \{0.08, 0.04, 0.00\}$.
È il confronto più diretto per quantificare il compromesso tra efficienza ed
equità dichiarato come obiettivo di questo lavoro.
Variando solo questo iperparametro della ricompensa, senza toccare stato o architettura,
osserveremo quanto costa in termini di travel time medio ottenere una data riduzione del travel time massimo e del max wait.

\subsection{Studio di ablazione}

Il secondo confronto (§\ref{sec:ablation-risultati}) isola sistematicamente ogni blocco di differenze
tra l'implementazione di questo lavoro e la procedura
descritta nell'articolo originale di MetaSTGAT, tramite 6 preset che
vanno dalla replica fedele del paper al modello MetaSTGAT Pro.
È un confronto necessario perché né l'uno né l'altro estremo, da soli,
spiegano il divario di prestazioni osservato tra i due: isolare un
sottosistema alla volta permette di attribuire quel divario a cause specifiche.

Ogni preset applica un gruppo coerente di scelte configurative,
isolando un sottosistema alla volta rispetto al modello MetaSTGAT Pro (§\ref{subsec:framing-pro}):

\begin{table}[htbp]
  \centering
  \small
  \begin{tabular}{@{}llp{8.5cm}@{}}
    \toprule
    Preset                   & Reward          & Cosa isola rispetto a Pro                                             \\
    \midrule
    \texttt{pro}             & \texttt{custom} & --                                                                    \\
    \texttt{temporal}        & \texttt{custom} & burn-in + BPTT sulla Meta-LSTM                                        \\
    \texttt{single\_dqn}     & \texttt{custom} & Double DQN                                                            \\
    \texttt{vanilla\_buffer} & \texttt{custom} & PER, Huber, soft update, grad clipping, warm-up, esplorazione ciclica \\
    \texttt{environment}     & \texttt{paper}  & campo visivo, feature dello stato $w_i^t$, maschera fasi consecutive  \\
    \texttt{paper}           & \texttt{paper}  & Replica fedele di Wang et al. (2022)                                  \\
    \bottomrule
  \end{tabular}
  \caption{Preset di ablazione e sottosistema isolato da ciascuno. La reward \texttt{custom} è definita da Eq. (\ref{eq:reward-custom}), mentre la reward \texttt{paper} è definita da Eq. (\ref{eq:reward-paper}).}
  \label{tab:ablation-preset}
\end{table}

\subsection{Studio sui meccanismi spaziali}

Il terzo studio (§\ref{sec:confronto-spaziale}) sostituisce il solo
meccanismo di comunicazione spaziale tra intersezioni di MetaSTGAT Pro.
Vengono testate le seguenti varianti già introdotte nel Capitolo \ref{sec:spatial-alt}:
\begin{itemize}
  \item Meta-GAT (il riferimento, a un layer come nella formulazione standard e a due layer)
  \item Meta-GCN (a uno e due layer)
  \item Meta-SONAR (a uno e due layer)
\end{itemize}

È il confronto che verifica se il beneficio osservato dipenda
specificamente dal meccanismo di attenzione appreso, o se una qualunque
forma di condivisione spaziale dell'informazione tra nodi vicini produca
risultati equivalenti.

\section{Risultati}
\label{sec:risultati-baseline}

I tre studi del precedente capitolo (§\ref{sec:confronti}) sono presentati
secondo uno schema comune: una tabella mediata sulle 3 configurazioni di
training, una corrispondente sulle 7 di test, un plot overlay che le
riassume visivamente, un plot della variazione percentuale delle metriche principali rispetto a
MaxPressure, gli eventuali casi particolari osservabili solo scendendo
alla singola configurazione, chiusi da una discussione che include, dove il risultato non
era quello atteso, una possibile spiegazione.
In queste tabelle, dallo studio sul peso $\alpha$ in poi, sono evidenziati
in grassetto i due risultati migliori di ogni colonna, Fixed-Time incluso;
nelle due tabelle del confronto di base (§\ref{sec:risultati-base}), che
hanno solo tre righe, resta invece in grassetto il solo valore migliore.

Fixed-Time, MaxPressure e MetaSTGAT Pro ($\alpha=0.08$, il modello di
riferimento di questo lavoro) compaiono in ciascuno dei tre studi.
Dunque, per non ripetere la stessa discussione su questi tre riferimenti di base,
la sottosezione seguente li confronta da soli,
mentre le sottosezioni successive approfondiranno i rispettivi confronti
(peso $\alpha$, ablazione, meccanismi spaziali).

\subsection{Fixed-Time, MaxPressure e Pro: il confronto di base}
\label{sec:risultati-base}

\begin{table}[htbp]
  \centering
  \begin{tabular}{@{}lrrrr@{}}
    \toprule
    Modello                       & $\overline{tt}$ (s) & $tt_{\max}$ (s) & wait$_{\max}$ (s) & $tc$ $(\%)$   \\
    \midrule
    Fixed-Time                    & 339.3               & 1030.0          & \textbf{470.0}    & 84.9          \\
    MaxPressure                   & \textbf{161.9}      & 1070.0          & 955.0             & \textbf{96.6} \\
    MetaSTGAT Pro ($\alpha=0.08$) & 176.7               & \textbf{675.0}  & 630.0             & 95.7          \\
    \bottomrule
  \end{tabular}
  \caption{Travel time medio, massimo, max wait e trip completion dei tre
    controllori presenti in ogni studio di questo capitolo, mediati sulle 3
    configurazioni di training.}
  \label{tab:base-train}
\end{table}

\begin{table}[htbp]
  \centering
  \begin{tabular}{@{}lrrrr@{}}
    \toprule
    Modello                       & $\overline{tt}$ (s) & $tt_{\max}$ (s) & wait$_{\max}$ (s) & $tc$ $(\%)$   \\
    \midrule
    Fixed-Time                    & 460.7               & 1496.1          & 1289.1            & 79.9          \\
    MaxPressure                   & \textbf{238.3}      & 1647.4          & 1626.0            & 88.5          \\
    MetaSTGAT Pro ($\alpha=0.08$) & 242.2               & \textbf{1049.6} & \textbf{954.0}    & \textbf{90.1} \\
    \bottomrule
  \end{tabular}
  \caption{Come Tabella~\ref{tab:base-train}, mediata sulle 7 configurazioni
    di test.}
  \label{tab:base-test}
\end{table}

Fixed-Time è il controllore peggiore su quasi ogni metrica di riferimento,
con l'unica eccezione del max wait in training (470.0s, il più basso dei tre):
la rotazione fissa delle fasi impone per costruzione un
limite alla durata di ogni attesa, un vincolo che le configurazioni di
training, meno dense, non arrivano a violare.
Questo approccio però non coordina il flusso tra le intersezioni, motivo per cui
questo risultato non si ripropone nella batteria di test dove il traffico è più intenso,
con il max wait che peggiora notevolmente.
Non osservando mai lo stato della rete, resta comunque il limite superiore di
travel time che qualunque strategia appresa dovrebbe poter migliorare.
MaxPressure e MetaSTGAT Pro lo battono nettamente sulle restanti metriche,
confermando che la regola reattiva alla base di MaxPressure coglie un segnale reale nello stato della rete.

Tra MaxPressure e MetaSTGAT Pro il confronto è più sfumato. In training e in tesing MaxPressure
resta il modello con il travel time medio più basso in assoluto (rispettivamente 161.9s e 238.3s),
coerente con la sua ottimalità di portata a livello di rete;
MetaSTGAT Pro paga per questo un travel time leggermente peggiore sia in training (176.7s, +9.1\%) sia in test (242.2s, +1.7\%).
Per quanto riguarda il max wait però il rapporto si inverte nettamente:
MetaSTGAT Pro registra un valore il 34\% più basso di MaxPressure in training e il 41\% più basso in test
(Tabelle~\ref{tab:base-train}-\ref{tab:base-test}).
Questi risultati introducono il compromesso al centro di questo lavoro, discusso nella sua forma
completa al variare di $\alpha$ nella prossima sottosezione.
Da notare è il fatto che rispetto alle prestazioni di MaxPressure, il modello MetaSTGAT Pro
ottiene risultati migliori nel training piuttosto che nel test: quanto ottenuto conferma sia
l'importanza di un controllo coordinato tra semafori quando la densità veicolare aumenta, sia
le capacità di generalizzazione acquisita dal modello MetaSTGAT Pro.

Il basso max wait di Fixed-Time in training merita un chiarimento: il suo
$tt_{\max}$ (1030.0s) resta più del doppio del suo stesso max wait
(470.0s), mentre per MaxPressure e per MetaSTGAT Pro i due valori restano
vicini (rispettivamente 1070.0s contro 955.0s, e 675.0s contro 630.0s). La
causa è strutturale: la rotazione fissa delle fasi non produce mai un
blocco estremo isolato, ma un'attesa moderata che si ripete, quasi
identica, a ogni incrocio della rete; il $tt_{\max}$ osservato è dunque la
somma di molti contributi distribuiti lungo l'intero percorso del
veicolo, non il riflesso di un singolo episodio. Negli altri due
controllori vale il caso opposto: un max wait elevato è quasi sempre un
episodio isolato su un solo incrocio, che da solo spiega quasi per intero
il $tt_{\max}$ corrispondente, con il resto del tragitto gestito in modo
efficiente.

Sul completamento ($tc$) la separazione più netta è quella tra Fixed-Time
e tutti gli altri modelli di questo capitolo: 84.9\%/79.9\%
(training/test) contro un intervallo compreso tra il 95.0\% e il 98.3\% in
training e tra l'88.9\% e il 95.8\% in test per ogni altra variante
(Tabella~\ref{tab:tc-adjusted}, a fine capitolo); l'unica
eccezione parziale è MetaSTGCN, il cui completamento più basso è discusso
a parte in §\ref{sec:confronto-spaziale}. Con questo unico scarto
rilevante già osservato qui, le sottosezioni seguenti non ripetono la
discussione di $tc$ configurazione per configurazione: si limitano a
segnalarne la tendenza generale, rimandando alla stessa tabella per il
dettaglio completo.

\subsection{Studio del peso \texorpdfstring{$\alpha$}{alpha}: risultati}
\label{sec:risultati-alpha}
\label{sec:equita-alpha}

Questa sottosezione confronta i due controllori di riferimento
(§\ref{sec:risultati-base}) con le tre varianti del modello Pro che
differiscono solo per il peso $\alpha$ della penalità anti-starvation,
con $\alpha \in \{0.08, 0.04, 0.00\}$.

\begin{table}[htbp]
  \centering
  \begin{tabular}{@{}lrrr@{}}
    \toprule
    Modello                       & $\overline{tt}$ (s) & $tt_{\max}$ (s) & wait$_{\max}$ (s) \\
    \midrule
    Fixed-Time                    & 339.3               & 1030.0          & \textbf{470.0}    \\
    MaxPressure                   & \textbf{161.9}      & 1070.0          & 955.0             \\
    MetaSTGAT Pro ($\alpha=0.08$) & 176.7               & \textbf{675.0}  & \textbf{630.0}    \\
    MetaSTGAT Pro ($\alpha=0.04$) & \textbf{170.1}      & \textbf{750.0}  & 690.0             \\
    MetaSTGAT Pro ($\alpha=0.00$) & 178.6               & 925.0           & 885.0             \\
    \bottomrule
  \end{tabular}
  \caption{Travel time medio, massimo e max wait, mediati
    sulle 3 configurazioni di training.}
  \label{tab:risultati-train}
\end{table}

\begin{table}[htbp]
  \centering
  \begin{tabular}{@{}lrrr@{}}
    \toprule
    Modello                       & $\overline{tt}$ (s) & $tt_{\max}$ (s) & wait$_{\max}$ (s) \\
    \midrule
    Fixed-Time                    & 460.7               & 1496.1          & 1289.1            \\
    MaxPressure                   & \textbf{238.3}      & 1647.4          & 1626.0            \\
    MetaSTGAT Pro ($\alpha=0.08$) & 242.2               & \textbf{1049.6} & \textbf{954.0}    \\
    MetaSTGAT Pro ($\alpha=0.04$) & \textbf{236.4}      & \textbf{1181.6} & \textbf{1125.9}   \\
    MetaSTGAT Pro ($\alpha=0.00$) & 248.7               & 1569.9          & 1539.4            \\
    \bottomrule
  \end{tabular}
  \caption{Come Tabella~\ref{tab:risultati-train}, mediata sulle 7
    configurazioni di test.}
  \label{tab:risultati-test}
\end{table}

Sul travel time medio ciò che verrebbe naturale attendersi è che
al crescese del valore di $\alpha$ la metrica in questione peggiori, con
le due configurazioni prive di penalità anti-starvation, MaxPressure e $\alpha=0.00$, che risulterebbero essere le
migliori, in quanto nessuna delle due paga un compromesso con l'equità, quindi
nessuna dovrebbe peggiorare travel time medio in cambio di un vantaggio che non
persegue.
Mentre questa analisi preliminare su $\alpha$ si conferma per i modelli con $\alpha=0.08$ e $\alpha=0.04$,
$\alpha=0.00$ risulta essere la peggiore delle tre varianti Pro su
questa metrica in entrambi i regimi (178.6s in training contro 176.7s di
$\alpha=0.08$, 248.7s in test contro 242.2s). Una causa plausibile è che
con $\alpha=0$ il termine di ricompensa legato all'attesa scompare, ma il
vettore $w_i^t$ resta nello stato osservato dalla rete
(§\ref{sec:formalizzazione}); la rete continua quindi a vedere un segnale
che la ricompensa non penalizza più causando un disallineamento tra stato e
obiettivo che può disturbare l'addestramento.

Un ulteriore risultato da sottolineare è quello ottenuto con $\alpha=0.04$ sul travel time medio in test, che non solo risulta essere il migliore tra le varianti del modello Pro,
ma supera anche MaxPressure (236.4s contro 238.3s), confermando che una quantità di penalità anti-starvation non nulla ma sufficientemente contenuta non ostacola necessariamente la performance su travel time medio.
Il risultato regge sia sulla media assoluta sia sulla media delle variazioni
percentuali per singola configurazione (Figura~\ref{fig:diffpct-alpha}),
quindi non è l'effetto di poche configurazioni favorevoli: è un segnale
che la struttura coordina bene il traffico, con un margine che cresce
proprio sulle topologie più complesse e dense.

\begin{figure}[htbp]
  \centering
  \begin{subfigure}[t]{0.48\textwidth}
    \centering
    \includegraphics[width=\linewidth]{pro_alpha_train.png}
    \caption{Configurazioni di training}
  \end{subfigure}
  \hfill
  \begin{subfigure}[t]{0.48\textwidth}
    \centering
    \includegraphics[width=\linewidth]{pro_alpha_test.png}
    \caption{Configurazioni di test}
  \end{subfigure}
  \caption{Travel time medio (fascia scura), max wait (fascia intermedia) e
    travel time massimo (fascia chiara) per lo sweep su $\alpha$, stessi dati
    delle Tabelle~\ref{tab:risultati-train}-\ref{tab:risultati-test}.}
  \label{fig:overlay-pro-alpha}
\end{figure}

\begin{figure}[htbp]
  \centering
  \begin{subfigure}[t]{0.48\textwidth}
    \centering
    \includegraphics[width=\linewidth]{diff_pct_alpha_train.png}
    \caption{Configurazioni di training}
  \end{subfigure}
  \hfill
  \begin{subfigure}[t]{0.48\textwidth}
    \centering
    \includegraphics[width=\linewidth]{diff_pct_alpha_test.png}
    \caption{Configurazioni di test}
  \end{subfigure}
  \caption{Variazione percentuale di travel time medio e massimo rispetto a
    MaxPressure (riferimento a 0\%) per lo studio su $\alpha$, calcolata per
    singola configurazione e poi mediata.}
  \label{fig:diffpct-alpha}
\end{figure}

Su max wait e $tt_{\max}$ invece il quadro è quello atteso: entrambe le
metriche peggiorano in modo pressoché monotono al diminuire di $\alpha$,
cioè al ridursi del peso dato al caso di starvation più estremo. Il
riferimento $\alpha=0.08$ ottiene sia il $tt_{\max}$ più basso in entrambi
i regimi, sia il max wait più basso, il 34\% in meno
di MaxPressure in training e il 41\% in meno in test. In test $\alpha=0.04$
resta comunque nettamente sotto Fixed-Time su entrambe le metriche (max
wait 1125.9s e $tt_{\max}$ 1181.6s, contro 1289.1s e 1496.1s);
$\alpha=0.00$ invece pareggia o peggiora leggermente i risultati di Fixed-Time
(max wait 1539.4s contro 1289.1s, $tt_{\max}$ 1569.9s contro 1496.1s).

Nonostante il vantaggio di $\alpha=0.04$ sul travel time medio, il resto
di questo lavoro adotta $\alpha=0.08$ come riferimento standard per gli
studi di ablazione e dei meccanismi spaziali
(§\ref{sec:ablation-risultati}, §\ref{sec:confronto-spaziale}). La ragione
è la priorità dichiarata tra le metriche di questo lavoro:
$\alpha=0.04$ è sistematicamente meno incisivo di
$\alpha=0.08$ su max wait e $tt_{\max}$, le due metriche di caso peggiore
che questa tesi tiene in primo piano, a fronte di un guadagno sul solo
travel time medio, la meno critica delle tre. $\alpha=0.08$ resta quindi
la scelta più coerente con l'obiettivo del lavoro, anche al costo di
qualche secondo in più di travel time medio.

\subsection{Studio di ablazione: risultati}
\label{sec:ablation-risultati}

Questa sottosezione confronta il modello Pro ($\alpha=0.08$, riferimento)
con i cinque preset di ablation presenti nella Tabella~\ref{tab:ablation-preset}:
\texttt{paper} (replica fedele del riferimento originale), \texttt{environment}
(campo di visibilità, $w_i^t$ e maschera azioni rimossi),
\texttt{temporal} (burn-in e BPTT disattivati), \texttt{single\_dqn}
(Double DQN sostituita da DQN) e \texttt{vanilla\_buffer} (PER, Huber, soft
update, gradient clipping, warm-up ed esplorazione ciclica disattivati),
tutti allo stesso $\alpha=0.08$ del riferimento.

\begin{table}[htbp]
  \centering
  \begin{tabular}{@{}lrrr@{}}
    \toprule
    Modello                          & $\overline{tt}$ (s) & $tt_{\max}$ (s) & wait$_{\max}$ (s) \\
    \midrule
    Fixed-Time                       & 339.3               & 1030.0          & \textbf{470.0}    \\
    MaxPressure                      & \textbf{161.9}      & 1070.0          & 955.0             \\
    Pro ($\alpha=0.08$, riferimento) & 176.7               & 675.0           & 630.0             \\
    \texttt{paper}                   & 187.8               & 950.0           & 915.0             \\
    \texttt{environment}             & 179.1               & 1065.0          & 1020.0            \\
    \texttt{temporal}                & 184.8               & \textbf{645.0}  & 605.0             \\
    \texttt{single\_dqn}             & 174.6               & \textbf{595.0}  & \textbf{485.0}    \\
    \texttt{vanilla\_buffer}         & \textbf{171.5}      & 695.0           & 560.0             \\
    \bottomrule
  \end{tabular}
  \caption{Travel time medio, massimo e max wait sui preset
    di ablazione, mediati sulle 3 configurazioni di training.}
  \label{tab:ablation-train}
\end{table}

\begin{table}[htbp]
  \centering
  \begin{tabular}{@{}lrrr@{}}
    \toprule
    Modello                          & $\overline{tt}$ (s) & $tt_{\max}$ (s) & wait$_{\max}$ (s) \\
    \midrule
    Fixed-Time                       & 460.7               & 1496.1          & 1289.1            \\
    MaxPressure                      & \textbf{238.3}      & 1647.4          & 1626.0            \\
    Pro ($\alpha=0.08$, riferimento) & 242.2               & 1049.6          & 954.0             \\
    \texttt{paper}                   & 271.1               & 1527.4          & 1513.3            \\
    \texttt{environment}             & 252.0               & 1595.6          & 1564.7            \\
    \texttt{temporal}                & 256.6               & 1169.1          & 1029.9            \\
    \texttt{single\_dqn}             & \textbf{240.6}      & \textbf{986.1}  & \textbf{838.7}    \\
    \texttt{vanilla\_buffer}         & 240.9               & \textbf{1048.3} & \textbf{804.0}    \\
    \bottomrule
  \end{tabular}
  \caption{Come Tabella~\ref{tab:ablation-train}, mediata sulle 7
    configurazioni di test.}
  \label{tab:ablation-test}
\end{table}

Sul travel time medio in training \texttt{vanilla\_buffer},
\texttt{single\_dqn}, il modello Pro ed \texttt{environment} restano in
una fascia stretta (171.5-179.1s). La stessa quasi coincidenza tra
\texttt{vanilla\_buffer}, \texttt{single\_dqn} e Pro si ripete
in test (240.6-242.2s), ma \texttt{environment} se ne stacca visibilmente (252.0s).
Il distacco non è uniforme sulle
sette configurazioni di test: sulle due a 200m è più del doppio rispetto
alla media delle cinque a 100m (+7.5\% contro +2.9\%, calcolato per
singola configurazione). La causa più diretta e più probabile è la disattivazione del
cutoff di visibilità propria di questo preset. Infatti, in seguito a
un addestramento comune su sole griglie a 100m, il passaggio a griglie a 200m
produce con il preset in questione degli stati di valori mai osservati prima, in quanto
\texttt{environment} osserva l'intera corsia e dunque registra un conteggio veicoli mai incontrato in training.
Gli altri preset invece limitano questo scarto numerico grazie al ridotto campo visivo,
aiutando di fatto il modello a generalizzare su reti con spaziatura maggiore rispetto a quella del training.

\begin{figure}[htbp]
  \centering
  \begin{subfigure}[t]{0.48\textwidth}
    \centering
    \includegraphics[width=\linewidth]{ablation_train.png}
    \caption{Configurazioni di training}
  \end{subfigure}
  \hfill
  \begin{subfigure}[t]{0.48\textwidth}
    \centering
    \includegraphics[width=\linewidth]{ablation_test.png}
    \caption{Configurazioni di test}
  \end{subfigure}
  \caption{Travel time medio, max wait e travel time massimo sui preset di
    ablation, stessi dati delle
    Tabelle~\ref{tab:ablation-train}-\ref{tab:ablation-test}.}
  \label{fig:overlay-ablation}
\end{figure}

\begin{figure}[htbp]
  \centering
  \begin{subfigure}[t]{0.48\textwidth}
    \centering
    \includegraphics[width=\linewidth]{diff_pct_ablation_train.png}
    \caption{Configurazioni di training}
  \end{subfigure}
  \hfill
  \begin{subfigure}[t]{0.48\textwidth}
    \centering
    \includegraphics[width=\linewidth]{diff_pct_ablation_test.png}
    \caption{Configurazioni di test}
  \end{subfigure}
  \caption{Variazione percentuale di travel time medio e massimo rispetto a
    MaxPressure (riferimento a 0\%) sui preset di ablation, calcolata per
    singola configurazione/seed e poi mediata.}
  \label{fig:diffpct-ablation}
\end{figure}

Su $tt_{\max}$ e $wait_{max}$, \texttt{paper} ed \texttt{environment} restano
nettamente indietro in entrambi i regimi, coerentemente con l'unico
elemento che condividono e che differisce da ogni altro preset: l'assenza del
termine legato all'attesa nella ricompensa.

Sempre sulle stesse due metriche in questione invece, i tre preset che lasciano intatti stato, ricompensa e maschera
(\texttt{temporal}, \texttt{single\_dqn}, \texttt{vanilla\_buffer}),
registrano in training i migliori risultati in assoluto tra tutti i preset considerati, modello Pro incluso.

Tuttavia, per quanto riguarda i risultati in regime di test
il quadro si inverte per \texttt{temporal}, con Pro che registra $tt_{\max}$ e $wait_{max}$ migliori.
Un'analisi sulle singole configurazioni rivela che il crollo è sistematico
e cresce con la scala della griglia: il suo $tt_{\max}$ medio per
configurazione resta vicino al riferimento sulle tre configurazioni 4x4
(855-975s) ma sale a 1218-1524s su 5x5 e 6x6.
La causa probabilmente da ricercare nel fatto che \texttt{temporal} disattiva BPTT e burn-in: senza propagazione del gradiente
lungo la sequenza né un periodo di adattamento dello stato ricorrente
all'inizio di ogni episodio, il modulo temporale della Meta-LSTM ha meno
modo di generalizzare il proprio stato interno a topologie di scala
diversa da quella di training, mentre lo stesso stato interno resta
adeguato sulla scala 4x4 già vista. La curva di training offre un
indizio indipendente coerente con questa lettura: la loss di
\texttt{temporal} resta stabilmente più alta di quella del riferimento
per l'intero addestramento, un divario compatibile con
un limite strutturale imposto dall'assenza di BPTT e burn-in.

Il vantaggio di \texttt{single\_dqn} e \texttt{vanilla\_buffer} su
$tt_{\max}$ e max wait regge anche scendendo alla singola configurazione,
senza che questo lavoro riesca a indicarne con certezza la causa. Resta
aperta una domanda distinta, che uno studio dedicato dovrebbe verificare:
se la generalizzazione di questi due preset regga altrettanto bene su
spaziature stradali diverse da quella di training. Un indizio non
conclusivo in questo senso viene dal max wait: il vantaggio di entrambi i
preset sul modello Pro, ampio sulle configurazioni 5x5 e 6x6
mai viste in training, si riduce fino a scendere al di sotto del modello Pro sulle due configurazioni
a 200m. Non possiamo pertanto dire con certezza che questi preset funzionino altrettanto bene su spaziature stradali diverse da quella di training.

\subsection{Studio sui meccanismi spaziali: risultati}
\label{sec:confronto-spaziale}

Questa sottosezione confronta il modello Pro (MetaSTGAT, un layer)
con le sue varianti a meccanismo spaziale alternativo definite in
§\ref{sec:spatial-alt}: MetaSTGAT a due layer, MetaSTGCN a uno e due
layer, e MetaSTSONAR con $L=2$ e $L=4$ ricorrenze. Tutti i modelli
condividono l'architettura residua, gli iperparametri e procedura di training.


\begin{table}[htbp]
  \centering
  \begin{tabular}{@{}lrrr@{}}
    \toprule
    Modello                          & $\overline{tt}$ (s) & $tt_{\max}$ (s) & wait$_{\max}$ (s) \\
    \midrule
    Fixed-Time                       & 339.3               & 1030.0          & \textbf{470.0}    \\
    MaxPressure                      & \textbf{161.9}      & 1070.0          & 955.0             \\
    MetaSTGAT, 1L (Pro, riferimento) & 176.7               & 675.0           & 630.0             \\
    MetaSTGAT, 2L                    & 174.6               & \textbf{635.0}  & \textbf{510.0}    \\
    MetaSTGCN, 1L                    & 248.5               & 1340.0          & 1280.0            \\
    MetaSTGCN, 2L                    & 217.8               & 1180.0          & 1160.0            \\
    MetaSTSONAR, $L{=}2$             & 170.6               & \textbf{630.0}  & 535.0             \\
    MetaSTSONAR, $L{=}4$             & \textbf{170.2}      & 660.0           & 525.0             \\
    \bottomrule
  \end{tabular}
  \caption{Travel time medio, massimo e max wait sui
    meccanismi spaziali, mediati sulle 3 configurazioni di training.}
  \label{tab:architetture-train}
\end{table}

\begin{table}[htbp]
  \centering
  \begin{tabular}{@{}lrrr@{}}
    \toprule
    Modello                          & $\overline{tt}$ (s) & $tt_{\max}$ (s) & wait$_{\max}$ (s) \\
    \midrule
    Fixed-Time                       & 460.7               & 1496.1          & 1289.1            \\
    MaxPressure                      & 238.3               & 1647.4          & 1626.0            \\
    MetaSTGAT, 1L (Pro, riferimento) & 242.2               & \textbf{1049.6} & \textbf{954.0}    \\
    MetaSTGAT, 2L                    & 246.6               & 1128.0          & 990.0             \\
    MetaSTGCN, 1L                    & 333.8               & 1659.0          & 1631.6            \\
    MetaSTGCN, 2L                    & 288.5               & 1457.6          & 1428.9            \\
    MetaSTSONAR, $L{=}2$             & \textbf{237.9}      & \textbf{1038.4} & \textbf{862.3}    \\
    MetaSTSONAR, $L{=}4$             & \textbf{238.3}      & 1119.9          & 1047.4            \\
    \bottomrule
  \end{tabular}
  \caption{Come Tabella~\ref{tab:architetture-train}, mediata sulle 7
    configurazioni di test.}
  \label{tab:architetture-test}
\end{table}

Il divario più netto di questo confronto separa MetaSTGCN dagli altri due
meccanismi. In test, entrambe le varianti
MetaSTGCN registrano valori nettamente peggiori su ogni colonna (travel time medio
288.5-333.8s contro 237.9-246.6s di MetaSTGAT/MetaSTSONAR), con un
$tt_{\max}$ e un $wait_{\max}$ che restano vicini a quelli di MaxPressure pur
trattandosi di un modello addestrato con lo stesso obiettivo di equità del
riferimento. La differenza concettuale descritta in §\ref{subsec:meta-gcn},
un peso di aggregazione fissato dalla sola topologia contro un punteggio di
compatibilità appreso per coppia di nodi, si traduce quindi in un divario di
prestazioni ampio su questo compito.
Un secondo layer in MetaSTGCN
migliora la variante a singolo layer su tutte le metriche riportate in
entrambi i regimi ($\overline{tt}$ di training 248.5$\to$217.8s,
$tt_{\max}$ di test 1659.0$\to$1457.6s, max wait di test
1631.6$\to$1428.9s): la profondità aggiuntiva compensa qui, almeno in
parte, l'assenza di un meccanismo di compatibilità appreso, pur restando
lontana dalle prestazioni di MetaSTGAT e MetaSTSONAR.

MetaSTGAT e MetaSTSONAR hanno prestazioni pressochè pari tra loro su travel
time medio e massimo; sul max wait, però, è MetaSTSONAR con $L=2$ a ottenere il
risultato migliore dell'intero confronto (862.3s in test), davanti
sia al modello Pro di riferimento (954.0s) sia alla sua variante a
due layer (990.0s).

Tra le due profondità di MetaSTSONAR, $L=2$ supera $L=4$ anche
sul $tt_{\max}$ in entrambi i regimi (630.0s contro 660.0s in training,
1038.4s contro 1119.9s in test), senza il vantaggio teorico di campo
recettivo più ampio atteso da una profondità doppia.

Aggiungere un secondo layer a MetaSTGAT, a differenza di MetaSTGCN, non produce qui un
miglioramento netto: la variante a due layer resta vicina al riferimento
su travel time medio e massimo ma peggiora leggermente il max wait in
test (990.0s contro 954.0s).

Un caso particolare nelle configurazioni di
training (Tabella~\ref{tab:wait-full}): MetaSTGCN a un layer ottiene un
max wait medio di 1280s, peggiore di MaxPressure stesso (955s), pur
condividendo con il riferimento la stessa ricompensa orientata
all'equità. È il solo modello di tutto questo lavoro, tra quelli
addestrati con il termine anti-starvation pieno, a non riuscire a battere
la baseline reattiva su questa metrica nel proprio stesso regime di
addestramento, coerentemente con la spiegazione architetturale data sopra: un
peso di aggregazione fissato dalla topologia non ha modo di imparare a
dare priorità al segnale del vicino più in sofferenza, per quanto la
ricompensa lo incentivi a farlo.

\begin{figure}[htbp]
  \centering
  \begin{subfigure}[t]{0.48\textwidth}
    \centering
    \includegraphics[width=\linewidth]{architectures_train.png}
    \caption{Configurazioni di training}
  \end{subfigure}
  \hfill
  \begin{subfigure}[t]{0.48\textwidth}
    \centering
    \includegraphics[width=\linewidth]{architectures_test.png}
    \caption{Configurazioni di test}
  \end{subfigure}
  \caption{Travel time medio, max wait e travel time massimo sui meccanismi
    spaziali, stessi dati delle
    Tabelle~\ref{tab:architetture-train}-\ref{tab:architetture-test}.}
  \label{fig:overlay-architetture}
\end{figure}

\begin{figure}[htbp]
  \centering
  \begin{subfigure}[t]{0.48\textwidth}
    \centering
    \includegraphics[width=\linewidth]{diff_pct_architetture_train.png}
    \caption{Configurazioni di training}
  \end{subfigure}
  \hfill
  \begin{subfigure}[t]{0.48\textwidth}
    \centering
    \includegraphics[width=\linewidth]{diff_pct_architetture_test.png}
    \caption{Configurazioni di test}
  \end{subfigure}
  \caption{Variazione percentuale di travel time medio e massimo rispetto a
    MaxPressure (riferimento a 0\%) sui meccanismi spaziali, calcolata per
    singola configurazione/seed e poi mediata.}
  \label{fig:diffpct-architetture}
\end{figure}

\subsection{Tabelle dei risultati.}
Sono qui di seguito riportati i risultati completi di tutti i modelli rispetto ciascuna configurazioni testata.

\begin{sidewaystable}
  \centering
  \tiny
  \begin{tabular}{@{}lrrrrrrrrrr@{}}
    \toprule
    Modello                   & Train 1 & Train 2 & Train 3 & 4x4 Pk & 4x4-200 F & 4x4-200 Pk & 5x5 F & 5x5 Pk & 6x6 F & 6x6 Pk \\
    \midrule
    Fixed-Time                & 84.6    & 85.5    & 87.9    & 84.7   & 84.0      & 81.9       & 84.4  & 82.6   & 81.5  & 80.4   \\
    MaxPressure               & 98.0    & 98.3    & 97.1    & 92.8   & 93.5      & 92.0       & 92.0  & 91.2   & 90.7  & 89.2   \\
    Pro ($\alpha$=0.08)       & 96.4    & 97.5    & 96.8    & 93.3   & 95.3      & 92.8       & 93.5  & 93.1   & 92.9  & 92.8   \\
    Pro ($\alpha$=0.04)       & 97.2    & 97.4    & 97.2    & 94.0   & 95.8      & 93.6       & 93.3  & 92.5   & 92.4  & 92.0   \\
    Pro ($\alpha$=0.00)       & 96.2    & 96.9    & 96.8    & 92.9   & 93.7      & 92.2       & 93.2  & 92.3   & 92.2  & 91.1   \\
    MetaSTGAT 2L              & 96.5    & 97.4    & 97.8    & 93.7   & 94.5      & 92.6       & 93.7  & 92.2   & 93.1  & 92.3   \\
    Ablation: paper           & 95.0    & 96.5    & 96.3    & 91.6   & 93.5      & 90.2       & 91.4  & 90.8   & 90.7  & 88.9   \\
    Ablation: temporal        & 95.8    & 95.6    & 97.0    & 92.7   & 93.4      & 91.6       & 93.2  & 92.5   & 92.8  & 92.6   \\
    Ablation: single\_dqn     & 96.8    & 97.6    & 97.0    & 93.6   & 95.6      & 93.4       & 93.3  & 92.3   & 92.5  & 92.2   \\
    Ablation: environment     & 95.9    & 96.6    & 95.7    & 91.1   & 92.8      & 90.8       & 92.1  & 91.7   & 91.4  & 89.2   \\
    Ablation: vanilla\_buffer & 96.1    & 97.7    & 97.5    & 93.8   & 95.7      & 92.9       & 93.7  & 92.8   & 92.9  & 92.0   \\
    MetaSTGCN 1L              & 90.6    & 91.1    & 91.7    & 87.2   & 88.3      & 86.2       & 88.9  & 85.8   & 86.7  & 85.1   \\
    MetaSTGCN 2L              & 93.3    & 94.8    & 93.6    & 91.0   & 92.6      & 89.8       & 91.4  & 90.2   & 90.6  & 89.4   \\
    MetaSTSONAR L=2           & 97.3    & 97.7    & 97.5    & 93.4   & 95.5      & 93.1       & 93.9  & 92.6   & 93.3  & 92.0   \\
    MetaSTSONAR L=4           & 97.2    & 97.6    & 97.6    & 93.4   & 95.1      & 92.6       & 93.2  & 92.9   & 93.3  & 92.9   \\
    \bottomrule
  \end{tabular}
  \caption{Percentuale di trip completion $tc$ per ogni modello
    e ogni configurazione, corretta escludendo dal
    denominatore i veicoli strutturalmente impossibilitati ad arrivare in
    tempo. Nessun modello scende sotto l'80\%; a parte Fixed-Time (baseline non
    adattiva) l'unico caso ricorrente sotto il 90\% è MetaSTGCN 1L.}
  \label{tab:tc-adjusted}
\end{sidewaystable}

\begin{sidewaystable}
  \centering
  \tiny
  \begin{tabular}{@{}lrrrrrrrrrr@{}}
    \toprule
    Modello                   & Train 1 & Train 2 & Train 3 & 4x4 Pk                                     & 4x4-200 F                                  & 4x4-200 Pk                                 & 5x5 F                                       & 5x5 Pk                                     & 6x6 F                                      & 6x6 Pk                                     \\
    \midrule
    Fixed-Time                & 351.6   & 342.9   & 323.3   & 437.3 {\fontsize{4}{4}\selectfont$\pm$6.4} & 415.5 {\fontsize{4}{4}\selectfont$\pm$6.7} & 438.8 {\fontsize{4}{4}\selectfont$\pm$4.9} & 459.7 {\fontsize{4}{4}\selectfont$\pm$4.7}  & 485.0 {\fontsize{4}{4}\selectfont$\pm$5.1} & 484.8 {\fontsize{4}{4}\selectfont$\pm$6.2} & 503.7 {\fontsize{4}{4}\selectfont$\pm$4.6} \\
    MaxPressure               & 165.4   & 160.2   & 160.0   & 218.8 {\fontsize{4}{4}\selectfont$\pm$2.3} & 197.3 {\fontsize{4}{4}\selectfont$\pm$4.2} & 247.8 {\fontsize{4}{4}\selectfont$\pm$2.6} & 217.8 {\fontsize{4}{4}\selectfont$\pm$3.3}  & 258.6 {\fontsize{4}{4}\selectfont$\pm$4.9} & 244.7 {\fontsize{4}{4}\selectfont$\pm$5.5} & 283.3 {\fontsize{4}{4}\selectfont$\pm$1.4} \\
    Pro ($\alpha$=0.08)       & 180.2   & 176.3   & 173.5   & 230.9 {\fontsize{4}{4}\selectfont$\pm$2.6} & 199.5 {\fontsize{4}{4}\selectfont$\pm$2.4} & 252.0 {\fontsize{4}{4}\selectfont$\pm$3.3} & 224.8 {\fontsize{4}{4}\selectfont$\pm$3.1}  & 266.7 {\fontsize{4}{4}\selectfont$\pm$3.4} & 243.3 {\fontsize{4}{4}\selectfont$\pm$3.4} & 278.6 {\fontsize{4}{4}\selectfont$\pm$4.0} \\
    Pro ($\alpha$=0.04)       & 172.5   & 169.4   & 168.4   & 223.3 {\fontsize{4}{4}\selectfont$\pm$3.1} & 196.9 {\fontsize{4}{4}\selectfont$\pm$3.4} & 248.2 {\fontsize{4}{4}\selectfont$\pm$3.8} & 213.4 {\fontsize{4}{4}\selectfont$\pm$4.6}  & 258.2 {\fontsize{4}{4}\selectfont$\pm$3.0} & 237.1 {\fontsize{4}{4}\selectfont$\pm$5.0} & 277.5 {\fontsize{4}{4}\selectfont$\pm$4.5} \\
    Pro ($\alpha$=0.00)       & 183.5   & 180.8   & 171.5   & 237.5 {\fontsize{4}{4}\selectfont$\pm$3.8} & 206.1 {\fontsize{4}{4}\selectfont$\pm$3.5} & 257.8 {\fontsize{4}{4}\selectfont$\pm$2.6} & 230.9 {\fontsize{4}{4}\selectfont$\pm$4.5}  & 269.8 {\fontsize{4}{4}\selectfont$\pm$3.7} & 247.9 {\fontsize{4}{4}\selectfont$\pm$4.6} & 290.7 {\fontsize{4}{4}\selectfont$\pm$6.0} \\
    MetaSTGAT 2L              & 181.3   & 175.0   & 167.6   & 231.8 {\fontsize{4}{4}\selectfont$\pm$3.5} & 203.7 {\fontsize{4}{4}\selectfont$\pm$2.1} & 252.7 {\fontsize{4}{4}\selectfont$\pm$2.9} & 228.8 {\fontsize{4}{4}\selectfont$\pm$3.7}  & 275.4 {\fontsize{4}{4}\selectfont$\pm$5.5} & 245.3 {\fontsize{4}{4}\selectfont$\pm$4.0} & 289.5 {\fontsize{4}{4}\selectfont$\pm$3.7} \\
    Ablation: paper           & 194.1   & 182.9   & 186.5   & 253.9 {\fontsize{4}{4}\selectfont$\pm$4.6} & 217.7 {\fontsize{4}{4}\selectfont$\pm$3.6} & 287.3 {\fontsize{4}{4}\selectfont$\pm$4.8} & 245.3 {\fontsize{4}{4}\selectfont$\pm$11.4} & 294.5 {\fontsize{4}{4}\selectfont$\pm$5.2} & 275.8 {\fontsize{4}{4}\selectfont$\pm$2.6} & 318.6 {\fontsize{4}{4}\selectfont$\pm$3.5} \\
    Ablation: temporal        & 189.8   & 190.0   & 174.6   & 247.4 {\fontsize{4}{4}\selectfont$\pm$6.4} & 213.6 {\fontsize{4}{4}\selectfont$\pm$7.0} & 262.5 {\fontsize{4}{4}\selectfont$\pm$6.8} & 240.3 {\fontsize{4}{4}\selectfont$\pm$6.0}  & 282.1 {\fontsize{4}{4}\selectfont$\pm$4.7} & 256.6 {\fontsize{4}{4}\selectfont$\pm$3.6} & 293.8 {\fontsize{4}{4}\selectfont$\pm$4.1} \\
    Ablation: single\_dqn     & 178.4   & 173.4   & 172.1   & 227.3 {\fontsize{4}{4}\selectfont$\pm$1.8} & 198.8 {\fontsize{4}{4}\selectfont$\pm$3.0} & 249.2 {\fontsize{4}{4}\selectfont$\pm$3.8} & 221.8 {\fontsize{4}{4}\selectfont$\pm$2.2}  & 266.7 {\fontsize{4}{4}\selectfont$\pm$5.0} & 241.1 {\fontsize{4}{4}\selectfont$\pm$4.3} & 279.0 {\fontsize{4}{4}\selectfont$\pm$3.5} \\
    Ablation: environment     & 182.3   & 181.7   & 173.1   & 241.0 {\fontsize{4}{4}\selectfont$\pm$4.3} & 216.7 {\fontsize{4}{4}\selectfont$\pm$7.8} & 268.3 {\fontsize{4}{4}\selectfont$\pm$3.0} & 228.8 {\fontsize{4}{4}\selectfont$\pm$2.4}  & 268.7 {\fontsize{4}{4}\selectfont$\pm$4.9} & 249.8 {\fontsize{4}{4}\selectfont$\pm$8.1} & 292.2 {\fontsize{4}{4}\selectfont$\pm$9.9} \\
    Ablation: vanilla\_buffer & 177.9   & 170.2   & 166.5   & 230.3 {\fontsize{4}{4}\selectfont$\pm$2.9} & 202.2 {\fontsize{4}{4}\selectfont$\pm$4.1} & 252.4 {\fontsize{4}{4}\selectfont$\pm$4.2} & 220.1 {\fontsize{4}{4}\selectfont$\pm$4.5}  & 264.8 {\fontsize{4}{4}\selectfont$\pm$3.9} & 238.8 {\fontsize{4}{4}\selectfont$\pm$3.2} & 278.9 {\fontsize{4}{4}\selectfont$\pm$6.4} \\
    MetaSTGCN 1L              & 250.5   & 254.6   & 240.5   & 325.6 {\fontsize{4}{4}\selectfont$\pm$8.1} & 290.4 {\fontsize{4}{4}\selectfont$\pm$6.6} & 343.6 {\fontsize{4}{4}\selectfont$\pm$7.6} & 310.6 {\fontsize{4}{4}\selectfont$\pm$6.3}  & 357.8 {\fontsize{4}{4}\selectfont$\pm$13.0} & 331.5 {\fontsize{4}{4}\selectfont$\pm$9.5} & 376.9 {\fontsize{4}{4}\selectfont$\pm$3.4} \\
    MetaSTGCN 2L              & 225.4   & 215.1   & 212.7   & 280.1 {\fontsize{4}{4}\selectfont$\pm$5.0} & 242.7 {\fontsize{4}{4}\selectfont$\pm$4.4} & 299.3 {\fontsize{4}{4}\selectfont$\pm$3.8} & 266.7 {\fontsize{4}{4}\selectfont$\pm$5.5}  & 308.6 {\fontsize{4}{4}\selectfont$\pm$9.4} & 290.7 {\fontsize{4}{4}\selectfont$\pm$6.0} & 331.7 {\fontsize{4}{4}\selectfont$\pm$2.7} \\
    MetaSTSONAR L=2           & 174.1   & 170.1   & 167.5   & 225.3 {\fontsize{4}{4}\selectfont$\pm$5.0} & 196.4 {\fontsize{4}{4}\selectfont$\pm$3.4} & 246.6 {\fontsize{4}{4}\selectfont$\pm$3.1} & 219.5 {\fontsize{4}{4}\selectfont$\pm$4.1}  & 261.3 {\fontsize{4}{4}\selectfont$\pm$4.6} & 235.4 {\fontsize{4}{4}\selectfont$\pm$3.9} & 281.0 {\fontsize{4}{4}\selectfont$\pm$1.8} \\
    MetaSTSONAR L=4           & 171.9   & 171.1   & 167.7   & 224.8 {\fontsize{4}{4}\selectfont$\pm$4.4} & 198.3 {\fontsize{4}{4}\selectfont$\pm$2.4} & 249.2 {\fontsize{4}{4}\selectfont$\pm$3.2} & 221.5 {\fontsize{4}{4}\selectfont$\pm$3.2}  & 261.5 {\fontsize{4}{4}\selectfont$\pm$2.7} & 238.5 {\fontsize{4}{4}\selectfont$\pm$2.9} & 274.5 {\fontsize{4}{4}\selectfont$\pm$2.6} \\
    \bottomrule
  \end{tabular}
  \caption{Travel time medio $\overline{tt}$ (s) per ogni modello e ogni configurazione. Sulle configurazioni di test, media $\pm$ deviazione standard su 5 seed diversi.}
  \label{tab:ttmedio-full}
\end{sidewaystable}

\begin{sidewaystable}
  \centering
  \tiny
  \begin{tabular}{@{}lrrrrrrrrrr@{}}
    \toprule
    Modello                   & Train 1 & Train 2 & Train 3 & 4x4 Pk                                        & 4x4-200 F                                     & 4x4-200 Pk                                    & 5x5 F                                         & 5x5 Pk                                        & 6x6 F                                         & 6x6 Pk                                        \\
    \midrule
    Fixed-Time                & 1020.0  & 975.0   & 1095.0  & 1410.0 {\fontsize{4}{4}\selectfont$\pm$49.3}  & 1431.0 {\fontsize{4}{4}\selectfont$\pm$30.9}  & 1290.0 {\fontsize{4}{4}\selectfont$\pm$19.0}  & 1620.0 {\fontsize{4}{4}\selectfont$\pm$32.9}  & 1482.0 {\fontsize{4}{4}\selectfont$\pm$17.5}  & 1653.0 {\fontsize{4}{4}\selectfont$\pm$32.0}  & 1587.0 {\fontsize{4}{4}\selectfont$\pm$37.2}  \\
    MaxPressure               & 915.0   & 1155.0  & 1140.0  & 1194.0 {\fontsize{4}{4}\selectfont$\pm$20.3}  & 1692.0 {\fontsize{4}{4}\selectfont$\pm$133.6} & 1446.0 {\fontsize{4}{4}\selectfont$\pm$105.0} & 1800.0 {\fontsize{4}{4}\selectfont$\pm$0.0}   & 1800.0 {\fontsize{4}{4}\selectfont$\pm$0.0}   & 1800.0 {\fontsize{4}{4}\selectfont$\pm$0.0}   & 1800.0 {\fontsize{4}{4}\selectfont$\pm$0.0}   \\
    Pro ($\alpha$=0.08)       & 780.0   & 600.0   & 645.0   & 888.0 {\fontsize{4}{4}\selectfont$\pm$63.9}   & 786.0 {\fontsize{4}{4}\selectfont$\pm$53.3}   & 876.0 {\fontsize{4}{4}\selectfont$\pm$7.3}    & 1242.0 {\fontsize{4}{4}\selectfont$\pm$152.8} & 1143.0 {\fontsize{4}{4}\selectfont$\pm$44.9}  & 1269.0 {\fontsize{4}{4}\selectfont$\pm$36.2}  & 1197.0 {\fontsize{4}{4}\selectfont$\pm$6.0}   \\
    Pro ($\alpha$=0.04)       & 675.0   & 750.0   & 825.0   & 1077.0 {\fontsize{4}{4}\selectfont$\pm$116.3} & 858.0 {\fontsize{4}{4}\selectfont$\pm$84.5}   & 1074.0 {\fontsize{4}{4}\selectfont$\pm$79.7}  & 1494.0 {\fontsize{4}{4}\selectfont$\pm$275.7} & 1227.0 {\fontsize{4}{4}\selectfont$\pm$24.0}  & 1302.0 {\fontsize{4}{4}\selectfont$\pm$105.4} & 1257.0 {\fontsize{4}{4}\selectfont$\pm$51.4}  \\
    Pro ($\alpha$=0.00)       & 885.0   & 1050.0  & 840.0   & 1275.0 {\fontsize{4}{4}\selectfont$\pm$195.1} & 1470.0 {\fontsize{4}{4}\selectfont$\pm$190.9} & 1317.0 {\fontsize{4}{4}\selectfont$\pm$119.4} & 1779.0 {\fontsize{4}{4}\selectfont$\pm$42.0}  & 1674.0 {\fontsize{4}{4}\selectfont$\pm$154.6} & 1722.0 {\fontsize{4}{4}\selectfont$\pm$99.7}  & 1665.0 {\fontsize{4}{4}\selectfont$\pm$82.2}  \\
    MetaSTGAT 2L              & 615.0   & 660.0   & 630.0   & 933.0 {\fontsize{4}{4}\selectfont$\pm$54.8}   & 714.0 {\fontsize{4}{4}\selectfont$\pm$43.1}   & 867.0 {\fontsize{4}{4}\selectfont$\pm$101.5}  & 1302.0 {\fontsize{4}{4}\selectfont$\pm$93.6}  & 1248.0 {\fontsize{4}{4}\selectfont$\pm$58.8}  & 1476.0 {\fontsize{4}{4}\selectfont$\pm$61.2}  & 1242.0 {\fontsize{4}{4}\selectfont$\pm$43.9}  \\
    Ablation: paper           & 930.0   & 975.0   & 945.0   & 1212.0 {\fontsize{4}{4}\selectfont$\pm$142.7} & 1257.0 {\fontsize{4}{4}\selectfont$\pm$93.1}  & 1326.0 {\fontsize{4}{4}\selectfont$\pm$145.3} & 1800.0 {\fontsize{4}{4}\selectfont$\pm$0.0}   & 1794.0 {\fontsize{4}{4}\selectfont$\pm$12.0}  & 1800.0 {\fontsize{4}{4}\selectfont$\pm$0.0}   & 1686.0 {\fontsize{4}{4}\selectfont$\pm$152.0} \\
    Ablation: temporal        & 630.0   & 705.0   & 600.0   & 975.0 {\fontsize{4}{4}\selectfont$\pm$43.5}   & 855.0 {\fontsize{4}{4}\selectfont$\pm$65.7}   & 948.0 {\fontsize{4}{4}\selectfont$\pm$47.8}   & 1350.0 {\fontsize{4}{4}\selectfont$\pm$93.0}  & 1218.0 {\fontsize{4}{4}\selectfont$\pm$36.0}  & 1524.0 {\fontsize{4}{4}\selectfont$\pm$47.1}  & 1281.0 {\fontsize{4}{4}\selectfont$\pm$75.1}  \\
    Ablation: single\_dqn     & 585.0   & 570.0   & 630.0   & 822.0 {\fontsize{4}{4}\selectfont$\pm$41.8}   & 774.0 {\fontsize{4}{4}\selectfont$\pm$57.4}   & 894.0 {\fontsize{4}{4}\selectfont$\pm$70.7}   & 1062.0 {\fontsize{4}{4}\selectfont$\pm$131.6} & 1065.0 {\fontsize{4}{4}\selectfont$\pm$75.3}  & 1212.0 {\fontsize{4}{4}\selectfont$\pm$39.6}  & 1086.0 {\fontsize{4}{4}\selectfont$\pm$63.4}  \\
    Ablation: environment     & 1020.0  & 1065.0  & 1110.0  & 1341.0 {\fontsize{4}{4}\selectfont$\pm$188.7} & 1446.0 {\fontsize{4}{4}\selectfont$\pm$191.3} & 1416.0 {\fontsize{4}{4}\selectfont$\pm$143.1} & 1746.0 {\fontsize{4}{4}\selectfont$\pm$58.2}  & 1689.0 {\fontsize{4}{4}\selectfont$\pm$132.0} & 1743.0 {\fontsize{4}{4}\selectfont$\pm$47.8}  & 1671.0 {\fontsize{4}{4}\selectfont$\pm$107.6} \\
    Ablation: vanilla\_buffer & 660.0   & 675.0   & 750.0   & 954.0 {\fontsize{4}{4}\selectfont$\pm$38.7}   & 795.0 {\fontsize{4}{4}\selectfont$\pm$42.4}   & 936.0 {\fontsize{4}{4}\selectfont$\pm$43.1}   & 1161.0 {\fontsize{4}{4}\selectfont$\pm$102.4} & 1071.0 {\fontsize{4}{4}\selectfont$\pm$33.7}  & 1215.0 {\fontsize{4}{4}\selectfont$\pm$26.8}  & 1182.0 {\fontsize{4}{4}\selectfont$\pm$22.0}  \\
    MetaSTGCN 1L              & 1590.0  & 1335.0  & 1095.0  & 1689.0 {\fontsize{4}{4}\selectfont$\pm$131.3} & 1593.0 {\fontsize{4}{4}\selectfont$\pm$204.4} & 1569.0 {\fontsize{4}{4}\selectfont$\pm$106.3} & 1536.0 {\fontsize{4}{4}\selectfont$\pm$169.3}   & 1728.0 {\fontsize{4}{4}\selectfont$\pm$122.7}   & 1728.0 {\fontsize{4}{4}\selectfont$\pm$81.8}   & 1770.0 {\fontsize{4}{4}\selectfont$\pm$60.0}   \\
    MetaSTGCN 2L              & 1215.0  & 1125.0  & 1200.0  & 1290.0 {\fontsize{4}{4}\selectfont$\pm$89.0}   & 1344.0 {\fontsize{4}{4}\selectfont$\pm$244.0} & 1263.0 {\fontsize{4}{4}\selectfont$\pm$99.2}  & 1719.0 {\fontsize{4}{4}\selectfont$\pm$102.0}   & 1434.0 {\fontsize{4}{4}\selectfont$\pm$134.7}   & 1692.0 {\fontsize{4}{4}\selectfont$\pm$114.4}   & 1461.0 {\fontsize{4}{4}\selectfont$\pm$121.0}   \\
    MetaSTSONAR L=2           & 585.0   & 585.0   & 720.0   & 915.0 {\fontsize{4}{4}\selectfont$\pm$79.9}   & 735.0 {\fontsize{4}{4}\selectfont$\pm$52.0}   & 879.0 {\fontsize{4}{4}\selectfont$\pm$46.1}   & 1077.0 {\fontsize{4}{4}\selectfont$\pm$95.1}   & 1170.0 {\fontsize{4}{4}\selectfont$\pm$80.5}   & 1278.0 {\fontsize{4}{4}\selectfont$\pm$29.1}   & 1215.0 {\fontsize{4}{4}\selectfont$\pm$128.0}   \\
    MetaSTSONAR L=4           & 600.0   & 675.0   & 705.0   & 945.0 {\fontsize{4}{4}\selectfont$\pm$122.6}   & 774.0 {\fontsize{4}{4}\selectfont$\pm$66.8}   & 939.0 {\fontsize{4}{4}\selectfont$\pm$122.5}  & 1233.0 {\fontsize{4}{4}\selectfont$\pm$150.7}   & 1311.0 {\fontsize{4}{4}\selectfont$\pm$97.9}   & 1392.0 {\fontsize{4}{4}\selectfont$\pm$58.8}   & 1245.0 {\fontsize{4}{4}\selectfont$\pm$51.1}   \\
    \bottomrule
  \end{tabular}
  \caption{Travel time massimo $tt_{\max}$ (s) per ogni modello e ogni configurazione. Sulle configurazioni di test, media $\pm$ deviazione standard su 5 seed diversi.}
  \label{tab:ttmax-full}
\end{sidewaystable}

\begin{sidewaystable}
  \centering
  \tiny
  \begin{tabular}{@{}lrrrrrrrrrr@{}}
    \toprule
    Modello                   & Train 1 & Train 2 & Train 3 & 4x4 Pk                                        & 4x4-200 F                                     & 4x4-200 Pk                                    & 5x5 F                                         & 5x5 Pk                                        & 6x6 F                                         & 6x6 Pk                                        \\
    \midrule
    Fixed-Time                & 465.0   & 285.0   & 660.0   & 984.0 {\fontsize{4}{4}\selectfont$\pm$54.2}   & 1053.0 {\fontsize{4}{4}\selectfont$\pm$75.5}  & 1014.0 {\fontsize{4}{4}\selectfont$\pm$64.8}  & 1572.0 {\fontsize{4}{4}\selectfont$\pm$80.7}  & 1314.0 {\fontsize{4}{4}\selectfont$\pm$102.4} & 1620.0 {\fontsize{4}{4}\selectfont$\pm$65.0}  & 1467.0 {\fontsize{4}{4}\selectfont$\pm$83.5}  \\
    MaxPressure               & 795.0   & 1050.0  & 1020.0  & 1146.0 {\fontsize{4}{4}\selectfont$\pm$29.4}  & 1680.0 {\fontsize{4}{4}\selectfont$\pm$138.1} & 1419.0 {\fontsize{4}{4}\selectfont$\pm$98.4}  & 1788.0 {\fontsize{4}{4}\selectfont$\pm$24.0}  & 1779.0 {\fontsize{4}{4}\selectfont$\pm$42.0}  & 1785.0 {\fontsize{4}{4}\selectfont$\pm$30.0}  & 1785.0 {\fontsize{4}{4}\selectfont$\pm$30.0}  \\
    Pro ($\alpha$=0.08)       & 735.0   & 540.0   & 615.0   & 798.0 {\fontsize{4}{4}\selectfont$\pm$46.9}   & 606.0 {\fontsize{4}{4}\selectfont$\pm$71.4}   & 795.0 {\fontsize{4}{4}\selectfont$\pm$21.2}   & 1077.0 {\fontsize{4}{4}\selectfont$\pm$194.3} & 1074.0 {\fontsize{4}{4}\selectfont$\pm$93.2}  & 1137.0 {\fontsize{4}{4}\selectfont$\pm$106.2} & 1191.0 {\fontsize{4}{4}\selectfont$\pm$18.0}  \\
    Pro ($\alpha$=0.04)       & 645.0   & 690.0   & 735.0   & 1050.0 {\fontsize{4}{4}\selectfont$\pm$127.3} & 687.0 {\fontsize{4}{4}\selectfont$\pm$69.3}   & 1005.0 {\fontsize{4}{4}\selectfont$\pm$94.9}  & 1479.0 {\fontsize{4}{4}\selectfont$\pm$281.4} & 1182.0 {\fontsize{4}{4}\selectfont$\pm$59.5}  & 1257.0 {\fontsize{4}{4}\selectfont$\pm$88.2}  & 1221.0 {\fontsize{4}{4}\selectfont$\pm$43.1}  \\
    Pro ($\alpha$=0.00)       & 855.0   & 990.0   & 810.0   & 1254.0 {\fontsize{4}{4}\selectfont$\pm$204.0} & 1428.0 {\fontsize{4}{4}\selectfont$\pm$229.7} & 1284.0 {\fontsize{4}{4}\selectfont$\pm$122.4} & 1779.0 {\fontsize{4}{4}\selectfont$\pm$42.0}  & 1665.0 {\fontsize{4}{4}\selectfont$\pm$165.4} & 1713.0 {\fontsize{4}{4}\selectfont$\pm$100.6} & 1653.0 {\fontsize{4}{4}\selectfont$\pm$75.5}  \\
    MetaSTGAT 2L              & 495.0   & 555.0   & 480.0   & 876.0 {\fontsize{4}{4}\selectfont$\pm$45.1}   & 603.0 {\fontsize{4}{4}\selectfont$\pm$58.8}   & 819.0 {\fontsize{4}{4}\selectfont$\pm$79.7}   & 972.0 {\fontsize{4}{4}\selectfont$\pm$183.3}  & 1239.0 {\fontsize{4}{4}\selectfont$\pm$49.8}  & 1236.0 {\fontsize{4}{4}\selectfont$\pm$221.0} & 1185.0 {\fontsize{4}{4}\selectfont$\pm$26.8}  \\
    Ablation: paper           & 885.0   & 945.0   & 915.0   & 1071.0 {\fontsize{4}{4}\selectfont$\pm$109.7} & 1215.0 {\fontsize{4}{4}\selectfont$\pm$91.5}  & 1278.0 {\fontsize{4}{4}\selectfont$\pm$135.3} & 1800.0 {\fontsize{4}{4}\selectfont$\pm$0.0}   & 1788.0 {\fontsize{4}{4}\selectfont$\pm$24.0}  & 1800.0 {\fontsize{4}{4}\selectfont$\pm$0.0}   & 1641.0 {\fontsize{4}{4}\selectfont$\pm$171.4} \\
    Ablation: temporal        & 570.0   & 675.0   & 570.0   & 819.0 {\fontsize{4}{4}\selectfont$\pm$112.1}  & 789.0 {\fontsize{4}{4}\selectfont$\pm$87.3}   & 882.0 {\fontsize{4}{4}\selectfont$\pm$58.0}   & 1074.0 {\fontsize{4}{4}\selectfont$\pm$72.6}  & 1188.0 {\fontsize{4}{4}\selectfont$\pm$25.8}  & 1236.0 {\fontsize{4}{4}\selectfont$\pm$189.6} & 1221.0 {\fontsize{4}{4}\selectfont$\pm$35.0}  \\
    Ablation: single\_dqn     & 510.0   & 450.0   & 495.0   & 726.0 {\fontsize{4}{4}\selectfont$\pm$30.9}   & 606.0 {\fontsize{4}{4}\selectfont$\pm$58.9}   & 759.0 {\fontsize{4}{4}\selectfont$\pm$76.8}   & 855.0 {\fontsize{4}{4}\selectfont$\pm$171.6}  & 981.0 {\fontsize{4}{4}\selectfont$\pm$104.2}  & 948.0 {\fontsize{4}{4}\selectfont$\pm$124.9}  & 996.0 {\fontsize{4}{4}\selectfont$\pm$111.7}  \\
    Ablation: environment     & 990.0   & 990.0   & 1080.0  & 1335.0 {\fontsize{4}{4}\selectfont$\pm$195.3} & 1422.0 {\fontsize{4}{4}\selectfont$\pm$211.6} & 1401.0 {\fontsize{4}{4}\selectfont$\pm$149.6} & 1734.0 {\fontsize{4}{4}\selectfont$\pm$51.6}  & 1671.0 {\fontsize{4}{4}\selectfont$\pm$142.5} & 1728.0 {\fontsize{4}{4}\selectfont$\pm$53.2}  & 1662.0 {\fontsize{4}{4}\selectfont$\pm$101.9} \\
    Ablation: vanilla\_buffer & 585.0   & 570.0   & 525.0   & 753.0 {\fontsize{4}{4}\selectfont$\pm$85.1}   & 594.0 {\fontsize{4}{4}\selectfont$\pm$50.7}   & 777.0 {\fontsize{4}{4}\selectfont$\pm$59.5}   & 768.0 {\fontsize{4}{4}\selectfont$\pm$235.3}  & 837.0 {\fontsize{4}{4}\selectfont$\pm$17.5}   & 834.0 {\fontsize{4}{4}\selectfont$\pm$99.8}   & 1065.0 {\fontsize{4}{4}\selectfont$\pm$70.4}  \\
    MetaSTGCN 1L              & 1485.0  & 1290.0  & 1065.0  & 1647.0 {\fontsize{4}{4}\selectfont$\pm$147.7}   & 1542.0 {\fontsize{4}{4}\selectfont$\pm$237.6} & 1509.0 {\fontsize{4}{4}\selectfont$\pm$94.2}  & 1527.0 {\fontsize{4}{4}\selectfont$\pm$169.5}   & 1719.0 {\fontsize{4}{4}\selectfont$\pm$134.0}   & 1716.0 {\fontsize{4}{4}\selectfont$\pm$86.2}   & 1761.0 {\fontsize{4}{4}\selectfont$\pm$70.7}   \\
    MetaSTGCN 2L              & 1185.0  & 1125.0  & 1170.0  & 1263.0 {\fontsize{4}{4}\selectfont$\pm$96.0}   & 1308.0 {\fontsize{4}{4}\selectfont$\pm$238.2} & 1203.0 {\fontsize{4}{4}\selectfont$\pm$85.1}  & 1713.0 {\fontsize{4}{4}\selectfont$\pm$97.9}   & 1401.0 {\fontsize{4}{4}\selectfont$\pm$141.2}   & 1674.0 {\fontsize{4}{4}\selectfont$\pm$118.7}   & 1440.0 {\fontsize{4}{4}\selectfont$\pm$131.5}   \\
    MetaSTSONAR L=2           & 510.0   & 510.0   & 585.0   & 753.0 {\fontsize{4}{4}\selectfont$\pm$146.5}    & 603.0 {\fontsize{4}{4}\selectfont$\pm$32.0}   & 720.0 {\fontsize{4}{4}\selectfont$\pm$9.5}    & 813.0 {\fontsize{4}{4}\selectfont$\pm$55.6}    & 1125.0 {\fontsize{4}{4}\selectfont$\pm$100.8}   & 831.0 {\fontsize{4}{4}\selectfont$\pm$64.1}    & 1191.0 {\fontsize{4}{4}\selectfont$\pm$149.3}   \\
    MetaSTSONAR L=4           & 555.0   & 525.0   & 495.0   & 882.0 {\fontsize{4}{4}\selectfont$\pm$134.0}   & 642.0 {\fontsize{4}{4}\selectfont$\pm$52.3}   & 816.0 {\fontsize{4}{4}\selectfont$\pm$125.4}  & 1137.0 {\fontsize{4}{4}\selectfont$\pm$187.4}   & 1263.0 {\fontsize{4}{4}\selectfont$\pm$83.5}   & 1368.0 {\fontsize{4}{4}\selectfont$\pm$84.5}   & 1224.0 {\fontsize{4}{4}\selectfont$\pm$43.1}   \\
    \bottomrule
  \end{tabular}
  \caption{Max wait $\text{wait}_{\max}$ (s) per ogni modello e ogni configurazione. Sulle configurazioni di test, media $\pm$ deviazione standard su 5 seed diversi. Sulle config \texttt{4x4-200}, valori a piena corsia.}
  \label{tab:wait-full}
\end{sidewaystable}
